\documentclass[10pt,a4paper]{amsart}
\usepackage{amsmath,amssymb,amsthm,graphicx}
% For commutative diagrams we use Xy-pic
\usepackage[all]{xy}

% We use 2cell for 2-commutative diagrams.
\xyoption{2cell}
% Theorem environments.
%
\theoremstyle{plain}
\newtheorem{theorem}[subsection]{Theorem}
\newtheorem{proposition}[subsection]{Proposition}
\newtheorem{lemma}[subsection]{Lemma}

\theoremstyle{definition}
\newtheorem{definition}[subsection]{Definition}
\newtheorem{example}[subsection]{Example}
\newtheorem{exercise}[subsection]{Exercise}
\newtheorem{situation}[subsection]{Situation}

\theoremstyle{remark}
\newtheorem{remark}[subsection]{Remark}
\newtheorem{remarks}[subsection]{Remarks}

\numberwithin{equation}{subsection}

% Macros
%
\def\lim{\mathop{\mathrm{lim}}\nolimits}
\def\colim{\mathop{\mathrm{colim}}\nolimits}
\def\Spec{\mathop{\mathrm{Spec}}}
\def\Hom{\mathop{\mathrm{Hom}}\nolimits}
\def\Ext{\mathop{\mathrm{Ext}}\nolimits}
\def\SheafHom{\mathop{\mathcal{H}\!\mathit{om}}\nolimits}
\def\SheafExt{\mathop{\mathcal{E}\!\mathit{xt}}\nolimits}
\def\Sch{\mathit{Sch}}
\def\Mor{\mathop{\mathrm{Mor}}\nolimits}
\def\Ob{\mathop{\mathrm{Ob}}\nolimits}
\def\Sh{\mathop{\mathit{Sh}}\nolimits}
\def\NL{\mathop{N\!L}\nolimits}
\def\CH{\mathop{\mathrm{CH}}\nolimits}
\def\proetale{{pro\text{-}\acute{e}tale}}
\def\etale{{\acute{e}tale}}
\def\QCoh{\mathit{QCoh}}
\def\Ker{\mathop{\mathrm{Ker}}}
\def\Im{\mathop{\mathrm{Im}}}
\def\Coker{\mathop{\mathrm{Coker}}}
\def\Coim{\mathop{\mathrm{Coim}}}

% Boxtimes
%
\DeclareMathSymbol{\boxtimes}{\mathbin}{AMSa}{"02}

%
% Macros for moduli stacks/spaces
%
\def\QCohstack{\mathcal{QC}\!\mathit{oh}}
\def\Cohstack{\mathcal{C}\!\mathit{oh}}
\def\Spacesstack{\mathcal{S}\!\mathit{paces}}
\def\Quotfunctor{\mathrm{Quot}}
\def\Hilbfunctor{\mathrm{Hilb}}
\def\Curvesstack{\mathcal{C}\!\mathit{urves}}
\def\Polarizedstack{\mathcal{P}\!\mathit{olarized}}
\def\Complexesstack{\mathcal{C}\!\mathit{omplexes}}
% \Pic is the operator that assigns to X its picard group, usage \Pic(X)
% \Picardstack_{X/B} denotes the Picard stack of X over B
% \Picardfunctor_{X/B} denotes the Picard functor of X over B
\def\Pic{\mathop{\mathrm{Pic}}\nolimits}
\def\Picardstack{\mathcal{P}\!\mathit{ic}}
\def\Picardfunctor{\mathrm{Pic}}
\def\Deformationcategory{\mathcal{D}\!\mathit{ef}}
\usepackage{cleveref}
\crefname{originalreference}{original reference}{original references}
\Crefname{originalreference}{Original reference}{Original references}
\makeatletter
\newlabel{lemma-characterize-finite-locally-constant}{{03RV}{source}{}{sourcepointer-1}{}}
\newlabel{lemma-characterize-finite-locally-constant@cref}{{[originalreference][0][]03RV}{[1][0][]source}{}{}{}}
\makeatother
\crefname{originalreference}{original reference}{original references}
\Crefname{originalreference}{Original reference}{Original references}
\makeatletter
\newlabel{lemma-colimit}{{00D5}{source}{}{sourcepointer-2}{}}
\newlabel{lemma-colimit@cref}{{[originalreference][0][]00D5}{[1][0][]source}{}{}{}}
\makeatother
\crefname{originalreference}{original reference}{original references}
\Crefname{originalreference}{Original reference}{Original references}
\makeatletter
\newlabel{lemma-directed-colimit-cohomology}{{03Q6}{source}{}{sourcepointer-3}{}}
\newlabel{lemma-directed-colimit-cohomology@cref}{{[originalreference][0][]03Q6}{[1][0][]source}{}{}{}}
\makeatother
\crefname{originalreference}{original reference}{original references}
\Crefname{originalreference}{Original reference}{Original references}
\makeatletter
\newlabel{lemma-equivalence-abelian-sheaves-point}{{04JQ}{source}{}{sourcepointer-4}{}}
\newlabel{lemma-equivalence-abelian-sheaves-point@cref}{{[originalreference][0][]04JQ}{[1][0][]source}{}{}{}}
\makeatother
\crefname{originallemma}{original lemma}{original lemmas}
\Crefname{originallemma}{Original lemma}{Original lemmas}
\makeatletter
\newlabel{limits-lemma-descend-etale}{{07RP}{source}{}{sourcepointer-5}{}}
\newlabel{limits-lemma-descend-etale@cref}{{[originallemma][0][]07RP}{[1][0][]source}{}{}{}}
\makeatother
\crefname{originallemma}{original lemma}{original lemmas}
\Crefname{originallemma}{Original lemma}{Original lemmas}
\makeatletter
\newlabel{limits-lemma-descend-finite-finite-presentation}{{01ZO}{source}{}{sourcepointer-6}{}}
\newlabel{limits-lemma-descend-finite-finite-presentation@cref}{{[originallemma][0][]01ZO}{[1][0][]source}{}{}{}}
\makeatother
\crefname{originallemma}{original lemma}{original lemmas}
\Crefname{originallemma}{Original lemma}{Original lemmas}
\makeatletter
\newlabel{limits-lemma-descend-finite-presentation}{{01ZM}{source}{}{sourcepointer-7}{}}
\newlabel{limits-lemma-descend-finite-presentation@cref}{{[originallemma][0][]01ZM}{[1][0][]source}{}{}{}}
\makeatother
\crefname{originalsection}{original section}{original sections}
\Crefname{originalsection}{Original section}{Original sections}
\makeatletter
\newlabel{pione-section-finite-etale-under-galois}{{03SF}{source}{}{sourcepointer-8}{}}
\newlabel{pione-section-finite-etale-under-galois@cref}{{[originalsection][0][]03SF}{[1][0][]source}{}{}{}}
\makeatother
\crefname{originallemma}{original lemma}{original lemmas}
\Crefname{originallemma}{Original lemma}{Original lemmas}
\makeatletter
\newlabel{sites-cohomology-lemma-torsors-h1}{{03AJ}{source}{}{sourcepointer-9}{}}
\newlabel{sites-cohomology-lemma-torsors-h1@cref}{{[originallemma][0][]03AJ}{[1][0][]source}{}{}{}}
\makeatother
\crefname{originalreference}{original reference}{original references}
\Crefname{originalreference}{Original reference}{Original references}
\makeatletter
\newlabel{theorem-vanishing-affine-curves}{{03SC}{source}{}{sourcepointer-10}{}}
\newlabel{theorem-vanishing-affine-curves@cref}{{[originalreference][0][]03SC}{[1][0][]source}{}{}{}}
\makeatother
\title{A cohomology class on a smooth affine curve over a field dies after a finite extension and a finite etale Galois cover of order dividing a power of the coefficient order}
\author{The Stacks Project Authors}
\date{stacks-project:etale-cohomology.tex:local-snapshot-20261011; GNU FDL 1.2 or later}
\begin{document}
\maketitle
\noindent\textit{Editorial note.} All delivered mathematical text is the author's own, copied from the pinned Stacks Project chapter file etale-cohomology.tex; no statement or proof step was rewritten, reordered or removed.\par\medskip
\noindent\textit{Editorial note.} The counted proof is the author's complete argument: the degree-above-one argument via the vanishing theorem and the passage to a finite subextension, the reduction to constant coefficients Z/dZ with trivial Galois action, the construction of the Galois cover from the class through the torsor dictionary, and the descent of that cover to a finite extension of k.\par\medskip
\noindent\textit{Editorial note.} The label lemma-smooth-base-change-fields names the counted statement, and the proof environment immediately following it in the source is delivered verbatim as its proof.\par\medskip
\noindent\textit{Editorial note.} The counted statement is carried by the author's own lemma environment in the chapter source.\par\medskip
\noindent\textit{Editorial note.} Pre-open gate: python operations/agent-authoring/tools/precard\_gate.py --source etale-cohomology.tex --label lemma-smooth-base-change-fields --start 17356 --end 17443 --local-id 3569 returned RESULT=OK (L3569 free, occupancy FREE, public interval clear). Candidate taken from lane-s4/halfB-priority.json.\par\medskip
\noindent\textit{Editorial note.} The proof contains the parenthetical '(small detail omitted)' at the point where the Galois group of the connected component is identified with a subgroup of Z/dZ; that verification is a formal group-theoretic check and the mathematical steps of the proof are written out, so the statement is counted, and the parenthetical is recorded here verbatim.\par\medskip
\noindent\textit{Editorial note.} The statement and proof are copied from the pinned chapter file; the cited inputs are kept as external pointers with their canonical Stacks tags.\par\medskip
\noindent\textit{Allowed standard prerequisites (editorial).}\begin{itemize}
\item Vanishing of the cohomology of an affine curve over a separably closed field. Same chapter, the theorem that the cohomology of a torsion coefficient sheaf on an affine curve over an algebraically closed field vanishes in degrees above one and is divisible in degree one. Supplies the vanishing of the pullback of the class to the algebraic closure, both in degree above one and, in degree one, the identification of the group with the corresponding group over the residue field.
\item Cohomology of a directed colimit and passage to a finite subextension. Same chapter, the lemma that a cohomology class defined over a limit of fields is defined over a finite subextension. Supplies the passage from the vanishing over the algebraic closure to the vanishing over a finite extension K' of K.
\item Equivalence between abelian sheaves on the spectrum of a field and discrete Galois modules. Same chapter, the lemma identifying sheaves on the etale site of a field with discrete modules over the Galois group. Supplies the description of the coefficient sheaf as a discrete module, which allows the reduction to finite n-torsion coefficients and to trivial Galois action after a finite extension.
\item Classification of torsors by the first cohomology group and representability of finite locally constant sheaves. Cohomology on Sites, the lemma that the first cohomology group with coefficients in a constant finite group classifies torsors; Same chapter, the lemma representing a finite locally constant sheaf by a finite etale scheme. Supply the construction of the finite etale morphism T -> X over the algebraic closure from the class, with its transitive action of Z/dZ on the fibre over the rational point.
\item Descent of finitely presented and etale objects along a limit of fields. Limits, the lemmas descending finite presentation, finite morphisms and etale morphisms along a field extension. Supplies the descent of the finite etale Galois cover from the algebraic closure to a finite extension k' of k.
\end{itemize}\medskip
\ifdefined\hypertarget\hypertarget{sourcepointer-1}{}\fi
\noindent\textit{Source reference.} Original reference 03RV: sources/stacks-etale-cohomology/pinned-original-source.tex, source lines 9994--9994. Prerequisite pointer invoked at proof line(s) 17419; it names a standard fact of the same chapter that the counted argument uses, and every step of the counted proof built on it is retained verbatim in the delivered body.\par\medskip
\ifdefined\hypertarget\hypertarget{sourcepointer-2}{}\fi
\noindent\textit{Source reference.} Original reference 00D5: sources/stacks-etale-cohomology/pinned-original-source.tex, source lines 7427--7427. Prerequisite pointer invoked at proof line(s) 17391; it names a standard fact of the same chapter that the counted argument uses, and every step of the counted proof built on it is retained verbatim in the delivered body.\par\medskip
\ifdefined\hypertarget\hypertarget{sourcepointer-3}{}\fi
\noindent\textit{Source reference.} Original reference 03Q6: sources/stacks-etale-cohomology/pinned-original-source.tex, source lines 7464--7464. Prerequisite pointer invoked at proof line(s) 17379; it names a standard fact of the same chapter that the counted argument uses, and every step of the counted proof built on it is retained verbatim in the delivered body.\par\medskip
\ifdefined\hypertarget\hypertarget{sourcepointer-4}{}\fi
\noindent\textit{Source reference.} Original reference 04JQ: sources/stacks-etale-cohomology/pinned-original-source.tex, source lines 9078--9078. Prerequisite pointer invoked at proof line(s) 17387; it names a standard fact of the same chapter that the counted argument uses, and every step of the counted proof built on it is retained verbatim in the delivered body.\par\medskip
\ifdefined\hypertarget\hypertarget{sourcepointer-5}{}\fi
\noindent\textit{Source reference.} Original lemma 07RP: sources/stacks-chow/tags.txt, source lines 9874--9874. Prerequisite pointer invoked at proof line(s) 17435; it names a standard fact of a sibling chapter that the counted argument uses, and every step of the counted proof built on it is retained verbatim in the delivered body.\par\medskip
\ifdefined\hypertarget\hypertarget{sourcepointer-6}{}\fi
\noindent\textit{Source reference.} Original lemma 01ZO: sources/stacks-chow/tags.txt, source lines 2601--2601. Prerequisite pointer invoked at proof line(s) 17434; it names a standard fact of a sibling chapter that the counted argument uses, and every step of the counted proof built on it is retained verbatim in the delivered body.\par\medskip
\ifdefined\hypertarget\hypertarget{sourcepointer-7}{}\fi
\noindent\textit{Source reference.} Original lemma 01ZM: sources/stacks-chow/tags.txt, source lines 2599--2599. Prerequisite pointer invoked at proof line(s) 17433; it names a standard fact of a sibling chapter that the counted argument uses, and every step of the counted proof built on it is retained verbatim in the delivered body.\par\medskip
\ifdefined\hypertarget\hypertarget{sourcepointer-8}{}\fi
\noindent\textit{Source reference.} Original section 03SF: sources/stacks-chow/tags.txt, source lines 4966--4966. Prerequisite pointer invoked at proof line(s) 17401; it names a standard fact of a sibling chapter that the counted argument uses, and every step of the counted proof built on it is retained verbatim in the delivered body.\par\medskip
\ifdefined\hypertarget\hypertarget{sourcepointer-9}{}\fi
\noindent\textit{Source reference.} Original lemma 03AJ: sources/stacks-chow/tags.txt, source lines 4312--4312. Prerequisite pointer invoked at proof line(s) 17414; it names a standard fact of a sibling chapter that the counted argument uses, and every step of the counted proof built on it is retained verbatim in the delivered body.\par\medskip
\ifdefined\hypertarget\hypertarget{sourcepointer-10}{}\fi
\noindent\textit{Source reference.} Original reference 03SC: sources/stacks-etale-cohomology/pinned-original-source.tex, source lines 15656--15656. Prerequisite pointer invoked at proof line(s) 17378; it names a standard fact of the same chapter that the counted argument uses, and every step of the counted proof built on it is retained verbatim in the delivered body.\par\medskip
\begin{lemma}
\label{lemma-smooth-base-change-fields}
Let $K/k$ be an extension of fields. Let $X$ be a smooth affine curve
over $k$ with a rational point $x \in X(k)$. Let $\mathcal{F}$ be an abelian
sheaf on $\Spec(K)$ annihilated by an integer $n$ invertible in $k$.
Let $q > 0$ and
$$
\xi \in H^q(X_K, (X_K \to \Spec(K))^{-1}\mathcal{F})
$$
There exist
\begin{enumerate}
\item finite extensions $K'/K$ and $k'/k$ with $k' \subset K'$,
\item a finite \'etale Galois cover $Z \to X_{k'}$ with group $G$
\end{enumerate}
such that the order of $G$ divides a power of $n$, such that
$Z \to X_{k'}$ is split over $x_{k'}$, and
such that $\xi$ dies in $H^q(Z_{K'}, (Z_{K'} \to \Spec(K))^{-1}\mathcal{F})$.
\end{lemma}

\begin{proof}
For $q > 1$ we know that $\xi$ dies in
$H^q(X_{\overline{K}}, (X_{\overline{K}} \to \Spec(K))^{-1}\mathcal{F})$
(Theorem \ref{theorem-vanishing-affine-curves}).
By Lemma \ref{lemma-directed-colimit-cohomology} we see that
this means there is a finite extension $K'/K$ such that
$\xi$ dies in $H^q(X_{K'}, (X_{K'} \to \Spec(K))^{-1}\mathcal{F})$.
Thus we can take $k' = k$ and $Z = X$ in this case.

\medskip\noindent
Assume $q = 1$. Recall that $\mathcal{F}$ corresponds to a
discrete module $M$ with continuous $\text{Gal}_K$-action, see
Lemma \ref{lemma-equivalence-abelian-sheaves-point}.
Since $M$ is $n$-torsion, it is the union of finite
$\text{Gal}_K$-stable subgroups. Thus
we reduce to the case where $M$ is a finite abelian group annihilated by $n$,
see Lemma \ref{lemma-colimit}. After replacing $K$ by a finite extension
we may assume that the action of $\text{Gal}_K$ on $M$ is trivial.
Thus we may assume $\mathcal{F} = \underline{M}$ is the constant
sheaf with value a finite abelian group $M$ annihilated by $n$.

\medskip\noindent
We can write $M$ as a direct sum of cyclic groups.
Any two finite \'etale Galois coverings whose Galois groups
have order invertible in $k$, can be dominated by a third one whose
Galois group has order invertible in $k$
(Fundamental Groups, Section \ref{pione-section-finite-etale-under-galois}).
Thus it suffices to prove the lemma when
$M = \mathbf{Z}/d\mathbf{Z}$ where $d | n$.

\medskip\noindent
Assume $M = \mathbf{Z}/d\mathbf{Z}$ where $d | n$.
In this case $\overline{\xi} = \xi|_{X_{\overline{K}}}$ is an element of
$$
H^1(X_{\overline{k}}, \mathbf{Z}/d\mathbf{Z}) =
H^1(X_{\overline{K}}, \mathbf{Z}/d\mathbf{Z})
$$
See Theorem \ref{theorem-vanishing-affine-curves}.
This group classifies $\mathbf{Z}/d\mathbf{Z}$-torsors, see
Cohomology on Sites, Lemma \ref{sites-cohomology-lemma-torsors-h1}.
The torsor corresponding to $\overline{\xi}$ (viewed as a sheaf on
$X_{\overline{k}, \etale}$) in turn gives rise to a finite \'etale morphism
$T \to X_{\overline{k}}$ endowed an action of $\mathbf{Z}/d\mathbf{Z}$
transitive on the fibre of $T$ over $x_{\overline{k}}$, see
Lemma \ref{lemma-characterize-finite-locally-constant}.
Choose a connected component $T' \subset T$ (if $\overline{\xi}$ has
order $d$, then $T$ is already connected).
Then $T' \to X_{\overline{k}}$ is a finite \'etale Galois
cover whose Galois group is a subgroup $G \subset \mathbf{Z}/d\mathbf{Z}$
(small detail omitted). Moreover the element $\overline{\xi}$ maps to zero
under the map $H^1(X_{\overline{k}}, \mathbf{Z}/d\mathbf{Z}) \to
H^1(T', \mathbf{Z}/d\mathbf{Z})$ as this is one
of the defining properties of $T$.

\medskip\noindent
Next, we use a limit argument to choose a finite extension $k'/k$
contained in $\overline{k}$ such that $T' \to X_{\overline{k}}$
descends to a finite \'etale Galois cover $Z \to X_{k'}$ with group $G$.
See Limits, Lemmas \ref{limits-lemma-descend-finite-presentation},
\ref{limits-lemma-descend-finite-finite-presentation}, and
\ref{limits-lemma-descend-etale}.
After increasing $k'$ we may assume that $Z$ splits over $x_{k'}$.
The image of $\xi$ in
$H^1(Z_{\overline{K}}, \mathbf{Z}/d\mathbf{Z})$ is zero by construction.
Thus by Lemma \ref{lemma-directed-colimit-cohomology}
we can find a finite subextension $\overline{K}/K'/K$
containing $k'$ such that $\xi$ dies in $H^1(Z_{K'}, \mathbf{Z}/d\mathbf{Z})$
and this finishes the proof.
\end{proof}

\end{document}
